% Verbatim copy of /Users/alingalatan/workspaces/grothe/handoff/manuscript/legacy/ch59-all-sector-audit-extension.tex (file dated 2026-10-25). Only this line was added.
% Consolidated source module (formerly Chapter 59). Owner: lower-bound/10-coupled-complements.tex
\section{All-sector anchors and a coupled-cubic extension}
\label{ch:all-sector-audit-extension}
\chapsummary{All-sector extensions on the stated cores and buffered tensor regions. Preserve the exceptional fresh mixed block and the branch-dependent retained deficits; there is no unrestricted even-sector removal.}

\paragraph{Provenance and status (IDE note, 21 October).} Lower-bound GPT all-sectors audit extension (\code{conjK-all-sector-audit-extension-20261021.zip}, archived in
\pathref{gpt-reviews/all-sector-audit-20261021/}). IDE replay \code{python -I verify\_all.py --full --jobs 6} (pinned \code{python-flint 0.9.0}, unrelated directory): ALL AUDIT AND
EXTENSION CHECKS PASS in 318\,s, with every incoming integral regenerated, all 46 one-driver anchors present, eight malformed-input rejections, the six arbitrary-cubic
all-sector regions, and the any-$N$ penalty $3/20$. IDE check of Lemma (Quadratic multiplication bounds): the sharp cases give exactly $17$ and $12$
($\E[|H|^2(|X|^2-1)^2]=17$, $\E[|H|^2|X|^4/2]=12$), and on 30 random coupled tensors on $\C^2$ (exact Gauss--Hermite moments) the largest ratios to the then-used $17e$, $e+10h$, $12e$
and $e+9h$ were $0.78$, $0.46$, $0.80$ and $0.37$. The canonical multiplication lemma below now replaces those influence coefficients.

\subsection{Scope and audit of the received result}
The input is \texttt{conjK-all-sectors-20261020.zip}. Its full source is
preserved in the accompanying audit directory. Statements labelled
\textbf{[P]} are analytic deductions; \textbf{[C]} denotes exact rational
checks using directed integral enclosures. These are not Lean results.

The assertion about all of $H_{2,1}(\mathbb C^2)$ means two coordinates
\emph{including} the visible coordinate: $(Z,X)$. It is not a theorem for
an arbitrary cubic on two remote coordinates together with $Z$.
The input proves, at the peak profile,
\[
 Q\in H_{2,1}(Z,X),\quad g\ge0
 \quad\Longrightarrow\quad \Lambda(Q+gU)\le S.
\]
Here $U$ is fresh. The six coefficients of $Q$ are unrestricted.
This strengthens the earlier small-buffer all-sector core theorem.
The earlier quantitative conditional-complement neighborhoods remain
useful and are not superseded by this membership statement.

The extension to arbitrary external even fields has an exception:
for one Wick driver, the block
\[
 \operatorname{span}\{Z\bar X V_i,\bar Z X V_i\}
\]
must be absent in the input theorem when
$c>19/100$, $g<41/250$, and $|\alpha|<2/5$.
This exception is absent on $(Z,X)$ alone. For two Wick drivers,
the input retains an even penalty $1/20$ on its anchor covers, not on
every branch of its global theorem. The trace branch gives only
$S-1/2000$; the large-buffer branch gives
$S-\|Q\|^2/2000$. The README and suggested insertion must carry these
qualifications, already present in its main theorem statement.

\paragraph{Checker qualification.}
The original one-driver scalar checker accepts two files containing only
\texttt{END 0}. This is a coverage-validation omission, not a
counterexample to the matrix bounds. The supplied audit adds a guard
requiring exactly the 37 compact and 9 prior anchors, no duplicates,
correct end counts and all matrix fields; it tests malformed inputs too.
The full replay and the guard are separate from the dependency theorems:
the incoming package does not reprove its older odd-sector covers.

\subsection{A common tangent permits simultaneous additions [P]}
Fix the peak parameters
\[
 a=699/1250,\quad b=179/1250,\quad p=2500/1891,
 \qquad L=Z(a+b|Z|^2).
\]
All Gaussian coordinates are proper with second moment one.
The visible $Z$ is independent of the remote Gaussian vector $X$;
$U$ is independent of both. Set
\[
 \begin{split}
 O&=q+(|Z|^2-1)r_A+Z^2r_B/\sqrt2,\\
 E&=ZY+\bar ZW+\nu H(Z),\qquad
 G=O+E+gU,
 \end{split}
\]
where $q\in H_{2,1}(X)$, $r_A\in H_{1,0}(X)$,
$r_B\in H_{0,1}(X)$, $Y\in H_{1,1}(X)$, and
$W\in H_{2,0}(X)$. All remote fields may depend on every same
coordinate; there is no mutual independence hypothesis.
Write $E_{AB}=\|r_A\|^2+\|r_B\|^2$ and
$E_{\rm ev}=\|Y\|^2+\|W\|^2+|\nu|^2$.

Put $\rho_g(u)=\mathbb E|\sqrt u+gU|$,
$K_g(u)=2\rho_g'(u)$ and $F_0=L+q$. Define the actual common tangent
\[
 \mathcal M_g(O;q)=\mathbb E\{2\rho_g(|F_0|^2)
       +K_g(|F_0|^2)(|L+O|^2-|F_0|^2)\}
       -p(\|O\|^2+g^2).
\]
Remote half-turn parity, followed by concavity of $\rho_g$, gives
\begin{equation}\label{eq:ase-parity}
 \Lambda(G)\le\mathcal M_g(O;q)
       +\mathbb E[K_g(|L+q|^2)|E|^2]-pE_{\rm ev}.
\end{equation}
Indeed the two smoothed arguments are $L+O+E$ and $L+O-E$.
Their squared norms sum to $2|L+O|^2+2|E|^2$.
This is an inequality for the \emph{common tangent}; a black-box
inequality only for $\Lambda(O+gU)$ cannot be substituted for it.
The proofs of the buffered coupled-tensor bounds do bound this same
$\mathcal M_g$, by bounding its contents and completing the A/B square.

\subsection{Multiplication by second chaos [P]}
\label{sec:quadratic-multiplication-refined}
Write $q=\sum T_{ij,k}{:}X_iX_j\bar X_k{:}$, with $T_{ij,k}=T_{ji,k}$,
and put $e=\|q\|^2=2\sum|T_{ij,k}|^2$ and
$\kappa=\E|q|^4-2e^2$. For the annihilators
$a_i=\partial_{X_i}$ and $b_i=\partial_{\bar X_i}$ set
\[
 d=(a_iq)_i,\quad f=(b_iq)_i,\quad
 P=\E[dd^*],\quad Q=\E[\bar f f^{\mathsf T}],\quad
 I=P+Q,\quad h=\|I\|_{\rm op},\quad
 p_*=\|P\|_{\rm op},\quad r_*=\|Q\|_{\rm op}.
\]
The conjugation in $Q$ puts both Grams in the same coordinate convention.
Thus $\operatorname{tr}P=2e$, $\operatorname{tr}Q=e$, $h\le3e$,
$p_*\le\min(h,2e)$, and $r_*\le\min(h,e)$.
The sharp influence theorem gives $\kappa\le14eh\le42e^2$.

\begin{lemma}[Refined quadratic multiplication bounds]
\label{lem:quadratic-multiplication-refined}
For every complex $Y\in H_{1,1}$ and $W\in H_{2,0}$ on the same
remote coordinates as $q$,
\begin{align}
 \E|qY|^2&\le (e+4h+2p_*)\|Y\|^2\le M_Y\|Y\|^2,
 &M_Y&=e+4h+2\min(h,2e),\label{eq:ase-MY}\\
 \E|qW|^2&\le (e+\tfrac{10}{3}h+r_*)\|W\|^2\le M_W\|W\|^2,
 &M_W&=e+\tfrac{10}{3}h+\min(h,e).\label{eq:ase-MW}
\end{align}
In particular $M_Y\le\min(17e,e+6h)$ and
$M_W\le\min(12e,e+13h/3)$. The variance constants $17,12$ are sharp;
no sharpness assertion is made for $6,13/3$.
No independence between $q,Y,W$ is assumed.
\end{lemma}
\begin{proof}
On the orthogonal complex Hermite basis,
$a_iH_{\mu,\nu}=\mu_iH_{\mu-e_i,\nu}$,
$a_i^*H_{\mu,\nu}=H_{\mu+e_i,\nu}$ and
$\|H_{\mu,\nu}\|^2=\mu!\nu!$, with the analogous formulas for $b_i$.
Multiplication by $X_i$ is $a_i^*+b_i$.

First let $Y$ be real. Its Hermitian coefficient matrix can be
unitarily diagonalized, so $Y=\sum_i\alpha_i(|X_i|^2-1)$ with real
$\alpha_i$ and $A^2=\sum\alpha_i^2=\|Y\|^2$.
The three orthogonal output grades of $qY$ are
\[
 Cq,\ Nq,\ Dq,\qquad
 C=\sum\alpha_i a_i^*b_i^*,\quad
 N=\sum\alpha_i(a_i^*a_i+b_i^*b_i),\quad
 D=\sum\alpha_i b_i a_i.
\]
Since $C^*=D$ and
$[D,C]=\sum\alpha_i^2(1+a_i^*a_i+b_i^*b_i)$,
\begin{equation}\label{eq:quadratic-Y-identity}
 \|qY\|^2=eA^2+\sum_i\alpha_i^2 I_{ii}+\|Nq\|^2+2\|Dq\|^2.
\end{equation}
On a degree-three basis vector write $m_i=\mu_i+\nu_i$.
Then $(\sum\alpha_i m_i)^2\le3\sum\alpha_i^2m_i$, hence
$\|Nq\|^2\le3\sum\alpha_i^2 I_{ii}\le3hA^2$.
The output coefficient of $X_\ell$ in $Dq$ is
$2\sum_i\alpha_iT_{i\ell,i}$. Contracting the output against a unit
vector $v$, put $S_{ik}=\sum_\ell\bar v_\ell T_{i\ell,k}$.
Then
\[
 |2\sum_i\alpha_i S_{ii}|^2\le4A^2\sum_{i,k}|S_{ik}|^2
 =A^2\|\sum_\ell\bar v_\ell a_\ell q\|^2\le p_*A^2.
\]
This proves the first bound for real $Y$. Apply it separately to
$\Re Y,\Im Y$ and add, using $|qY|^2=|q|^2((\Re Y)^2+(\Im Y)^2)$,
to cover every complex $Y$.

For $W$, unitary Takagi diagonalization gives
$W=\sum\beta_iX_i^2$, $\beta_i\ge0$, and
$\|W\|^2=2B^2$ with $B^2=\sum\beta_i^2$.
Its two output grades are $C_Bq$ and $2Mq$, where
\[
 C_B=\sum\beta_i(a_i^*)^2,\quad D_B=C_B^*=\sum\beta_i a_i^2,
 \quad M=\sum\beta_i a_i^*b_i.
\]
Put $x=\sum\beta_i^2\|a_iq\|^2$ and $z=\sum\beta_i^2\|b_iq\|^2$.
Commuting once gives
\[
 \|C_Bq\|^2=2eB^2+4x+\|D_Bq\|^2,\quad
 \|Mq\|^2=z+S,\quad
 S=\sum_{i,j}\beta_i\beta_j\langle a_jb_iq,a_ib_jq\rangle.
\]
The sum $S$ is real. Cauchy--Schwarz, assigning the weights to either
factor, gives both
\[
 S\le\sum_{i,j}\beta_i^2\|a_jb_iq\|^2=2z,
 \qquad S\le\sum_{i,j}\beta_i^2\|a_ib_jq\|^2=x.
\]
Therefore $S\le\min(x,2z)\le2(x+z)/3$ and
$x+z=\sum\beta_i^2 I_{ii}\le hB^2$.
The output coefficient of $\bar X_k$ in $D_Bq$ is
$2\sum_i\beta_iT_{ii,k}$. Dualizing the remaining $k$ index and
applying Cauchy--Schwarz over $(i,j)$ gives
$\|D_Bq\|^2\le2r_*B^2$. Consequently
\begin{equation}\label{eq:quadratic-W-identity}
 \|qW\|^2=2eB^2+4(x+z)+4S+\|D_Bq\|^2
 \le(2e+20h/3+2r_*)B^2.
\end{equation}
Divide by $2B^2$. For the variance sharpness take
$q=H(X)$, $Y=|X|^2-1$, $W=X^2/\sqrt2$.
Then $e=1$, $h=3$, $p_*=2$, $r_*=1$, and the two products have
squared norms $17$ and $12$ by the exponential moment table.
\end{proof}

\paragraph{Retaining the separate influence norms [P].}
The same proof gives the additional refinements
\begin{align*}
 M_Y^{\rm dir}&=e+h+\min\{3h,(\sqrt{2p_*}+\sqrt{r_*})^2\}+2p_*,\\
 M_W^{\rm dir}&=e+2h+2\min\{p_*,2r_*,2h/3\}+r_*.
\end{align*}
For $Y$, the holomorphic and antiholomorphic number-operator parts of
$Nq$ have norm at most $\sqrt{2p_*}A$ and $\sqrt{r_*}A$.
For $W$, retain $x\le p_*B^2$ and $z\le r_*B^2$ when bounding $S$.
These refinements help anisotropic tensors without changing the
scalar $(e,h)$ envelopes used in the certificates below.

\paragraph{Calibration and exact audit.}
For $q=N^{-1/2}\sum H(X_j)$ and the corresponding normalized sums
$Y=N^{-1/2}\sum(|X_j|^2-1)$, $W=(2N)^{-1/2}\sum X_j^2$,
\[
 e=1,\quad h=3/N,\quad \E|qY|^2=1+16/N,\quad
 \E|qW|^2=1+11/N.
\]
Thus any coefficients in universal bounds $e+C_Yh,e+C_Wh$ must satisfy
$C_Y\ge16/3$, $C_W\ge11/3$, even in the diffuse limit. We do not
claim these lower limits are universal upper bounds.
The portable checker verifies both exact operator identities on 240
Gaussian-integer Hermite polynomials in dimensions one through six,
and these sharpness formulas. These finite checks audit normalization;
the dimension-free proof is the operator argument above.
The refinement was supplied in \code{quadratic-multiplication.pdf};
this section replaces the former $e+10h,e+9h$ proof.

\subsection{A three-by-three majorant for every coupled cubic [P]}
Use the visible transforms and nonnegative spectral weight from the
received certificate:
\[
 K_g(|w|^2)=\int_0^\infty W_g(k)J_0(k|w|)\,dk,
 \qquad W_g(k)=e^{-x^2}-\sqrt\pi x\,\operatorname{erfc}(x),\quad x=gk/2.
\]
The visible transforms are
\[
 V_T=\mathbb E[TJ_0(k|L|)],\quad
 V_n=\mathbb E[T(T-2)^2J_0(k|L|)/2],\quad
 V_{2T}=\mathbb E[TJ_2(k|L|)],
\]
\[
 V_{Tn}=\mathbb E[T(T-2)J_0(k|L|)/\sqrt2],\quad
 V_{2Tn}=\mathbb E[T(T-2)J_2(k|L|)/\sqrt2].
\]
At a lower buffer $g_0>0$ let $K_T=\int W_{g_0}V_T$,
$K_n=\int W_{g_0}V_n$, and define
\[
 N_T=\int k^2W_{g_0}(V_T)_-,\quad
 N_n=\int k^2W_{g_0}(V_n)_-,\quad
 N_{2T}=\int k^2W_{g_0}|V_{2T}|,
\]
\[
 N_{Tn}=\int k^2W_{g_0}|V_{Tn}|,\qquad
 N_{2Tn}=\int k^2W_{g_0}|V_{2Tn}|.
\]
All integrals have domain $[0,\infty)$.
Define the real symmetric matrix
\[
 \mathcal Q=
 \begin{pmatrix}
 K_T+N_TM_Y/4&N_{2T}\sqrt{M_YM_W}/4&N_{Tn}\sqrt{e^2+\kappa}/4\\
 N_{2T}\sqrt{M_YM_W}/4&K_T+N_TM_W/4&N_{2Tn}\sqrt{2e^2+\kappa}/8\\
 N_{Tn}\sqrt{e^2+\kappa}/4&N_{2Tn}\sqrt{2e^2+\kappa}/8&K_n+N_ne/4
 \end{pmatrix}.
\]

\begin{theorem}[General even-sector extension]
For every $g\ge g_0$, and arbitrary shared remote coefficients,
\[
 \mathbb E[K_g(|L+q|^2)|E|^2]
 \le (\|Y\|,\|W\|,|\nu|)\mathcal Q
                  (\|Y\|,\|W\|,|\nu|)^\top.
\]
Consequently $\mathcal Q\preceq(p-\tau)I$ gives
\[
 \Lambda(O+E+gU)\le\mathcal M_g(O;q)-\tau E_{\rm ev}.
\]
\end{theorem}

\begin{proof}
First take $g=g_0$. The common remote phase has charges
1 for $q$, 0 for $Y$ and 2 for $W$. The diagonal Fourier
factors are $\mathbb E[\cos(k\Re q)|Y|^2]$ and its $W$
analogue. Their deviations from their unweighted values lie
between $-k^2\mathbb E|qY|^2/4$ and zero. This follows from
$1-\cos t\le t^2/2$ and phase averaging. Thus only the
negative parts of the visible transforms contribute to an
upper bound. The $\nu$ diagonal is the same argument with $e$.
For the $Y,W$ cross term the unweighted expectation vanishes;
Cauchy--Schwarz with the common nonnegative weight
$1-\cos(k\Re q)$ gives $k^2\sqrt{M_YM_W}/4$.

For the $Y,\nu$ cross term, phase averaging gives
$\mathbb E[YJ_0(k|q|)]$. The function
$u\mapsto J_0(k\sqrt u)$ is $k^2/4$-Lipschitz, including at
zero, since $|J_1(t)|\le |t|/2$. The independent-copy
variance identity implies
\[
 \operatorname{Var}(J_0(k|q|))
 \le k^4\operatorname{Var}(|q|^2)/16
 = k^4(e^2+\kappa)/16.
\]
Use $\mathbb EY=0$ and Cauchy--Schwarz.
For $W,\nu$, the charged phase average is a unit phase times
$J_2(k|q|)$. The integral representation of $J_2$ gives
$|J_2(t)|\le t^2/8$, so this term is bounded by
$k^2\sqrt{\mathbb E|q|^4}\|W\|/8$.
Integration against the corresponding absolute visible transforms
proves the matrix claim. These steps use no factorization of $q$.

Finally $K_g(u)$ decreases pointwise in $g$, by its positive
Laplace representation. Hence the whole nonnegative quadratic
form $\mathbb E[K_g|E|^2]$ is bounded by its value at $g_0$.
No sign assumption on an individual Hankel transform is used.
\end{proof}

\subsection{Certified consequences [P,C]}
The directed visible tables are freshly generated from the incoming
engine. Buffer endpoints are rounded \emph{down} before evaluation.
The checker \texttt{check\_general.py} uses exact fractions and
integer-square-root bounds, and checks all three positive pivots of
$(p-\tau)I-\mathcal Q_{\rm up}$. Entrywise upper bounds are used
only on the nonnegative norm vector, not asserted to be Loewner bounds.
Monotonicity in $e,h,\kappa$ covers the entire parameter rectangles.

The existing common-tangent A/B proofs therefore extend as follows.
In every row below, $q$ is an \emph{arbitrary} coupled remote cubic
in any dimension and $r_A,r_B,Y,W,\nu$ are unrestricted. The
conclusion is $\Lambda(G)<S+\delta-E_{AB}/2000-\tau E_{\rm ev}$:
\[
\begin{array}{c|c|c|c}
g\text{ at least}& e\text{ at most}&\delta&\tau\\\hline
37/200&9/4000&89/10^6&1/10\\
1/5&7/2500&89/10^6&1/10\\
9/40&49/10000&-1/200000&2/25\\
1/4&9/1000&-1/200000&1/40
\end{array}
\]
These use the variance-only coupled-tensor results (CT.5) as
common-tangent inputs. There are also the influence-only extensions
\[
 \boxed{g\ge9/40,\ e\le1/25,\ h\le9e/100
 \ \Longrightarrow\
 \Lambda(G)<S-1/12500-E_{AB}/2000-E_{\rm ev}/20,}
\]
\[
 \boxed{g\ge1/5,\ e\le1/25,\ h\le e/150
 \ \Longrightarrow\
 \Lambda(G)<S+89/10^6-E_{AB}/2000-9E_{\rm ev}/100.}
\]
The older certificates used $\kappa\le50eh/3$; the canonical bound
$\kappa\le14eh$ now gives ratios $63/50$ and $7/75$ instead.
No additional independence among the remote fields is introduced.

\begin{corollary}[Larger all-sector influence region and stronger penalties, \st{P,C}]
\label{cor:quadratic-all-sector-refinement}
For the same arbitrary coupled-cubic family, with all A/B/even/diagonal
amplitudes and phases unrestricted,
\[
 \begin{array}{c|c|c|c|c}
 g\text{ at least}&e\text{ at most}&h/e\text{ at most}&\delta&\tau\\\hline
 9/40&1/25&9/100&-1/12500&19/200\\
 9/40&1/25&1/10&-1/40000&9/100\\
 1/5&1/25&1/150&89/10^6&99/1000
 \end{array}
\]
Each row proves $\Lambda(G)<S+\delta-E_{AB}/2000-\tau E_{\rm ev}$.
In particular the first retained even penalty rises from $1/20$ to
$19/200$, and the middle row adds every even and diagonal sector to
Corollary~\ref{cor:sharp-influence-region}.
\end{corollary}
\begin{proof}
Insert the new $M_Y,M_W$ and $\kappa\le14eh$ into $\mathcal Q$.
The exact checker verifies all three positive LDL pivots of
$(p-\tau)I-\mathcal Q_{\rm up}$ at each endpoint; monotonicity covers
the whole parameter rectangle and every larger fresh buffer.
Combine with the existing common-tangent A/B bounds, including the
$h\le e/10$ completion of Corollary~\ref{cor:sharp-influence-region}.
The weak row retains that completion's remote-comparison dependency.
The four variance-only rows above are also replayed unchanged.

The old matrix certificates remain valid: $\mathcal Q$ has nonnegative
off-diagonal entries, and every changed entry decreases. For real $v$,
$v^\top\mathcal Q_{\rm new}v\le|v|^\top\mathcal Q_{\rm new}|v|
\le|v|^\top\mathcal Q_{\rm old}|v|$.
This is not an assertion that their difference is positive semidefinite.
\end{proof}
The exact rational replay and optional fresh visible-integral generator
are \code{verify\_quadratic.py} and \code{regenerate\_quadratic.py} in
\pathref{/Users/alingalatan/Documents/Codex/2026-10-10/plea/outputs/k-lower-audit-20261021/}.
They keep the directed quadrature engine and all odd-sector dependencies
explicit; they do not reprove the historical anchor covers.

\paragraph{A stronger penalty in the received any-$N$ theorem.}
For $\max c_j\le1/50$, $e\le1/100$ and
$7/200\le g^2\le9/200$, the input's same $3\times3$ Wick matrix
satisfies $\mathcal Q_{\rm Wick}\prec(p-3/20)I$.
The checker verifies this for each of its three variance thresholds.
Therefore its retained even penalty can be increased from $1/10$
to $3/20$, with every other coefficient and hypothesis unchanged.
This is a 50 percent improvement to that penalty, not to the
Grothendieck lower bound.

\subsection{What this changes in the endpoint strategy}
These are simultaneous all-sector statements; they avoid inserting an
even-sector comparison after an A/B baseline has become visible-dependent.
They remove that interaction obligation in the displayed coupled-tensor
regions as well as in the received Wick regions.
They do not remove the fresh-buffer hypothesis. An arbitrary finite
cubic need not contain such a buffer.

The actual-response constraints from the preceding work remain relevant:
at $t\ge1057/1250$, remote variance is less than $0.040726$ and
total cubic energy is less than $0.071981$. In particular a remote
Wick coefficient must be below $\sqrt{0.040726}<0.202$.
The downward region $\max c_j\ge0.315$ is therefore irrelevant to
actual responses at that peak, although it is a valid raw-K theorem.
The important open region remains a general shared cubic with small or
zero fresh buffer. The original-response support and projection cuts
address precisely that broader region and should remain the main plan.
No new endpoint follows until an exhaustive reduction and the complete
row cover have been supplied.
